% Unicode corpus (M3): polyglossia with Hebrew (right to left, bidi),
% Culmus David CLM by file name.
\documentclass{article}
\usepackage{fontspec}
\setmainfont{LibertinusSerif}[Extension=.otf, UprightFont=*-Regular,
  BoldFont=*-Bold, ItalicFont=*-Italic, BoldItalicFont=*-BoldItalic]
\newfontfamily\hebrewfont{DavidCLM}[Script=Hebrew, Extension=.otf,
  UprightFont=*-Medium, BoldFont=*-Bold, ItalicFont=*-MediumItalic,
  BoldItalicFont=*-BoldItalic]
\usepackage{polyglossia}
\setdefaultlanguage{english}
\setotherlanguage{hebrew}
\begin{document}
English first, then Hebrew inline: \texthebrew{שלום עולם}, and back.

\begin{hebrew}
זהו טקסט בעברית עם מספרים 123 ו\textbf{הדגשה}. משפט שני, ארוך יותר, כדי
שהפסקה תתפרס על פני יותר משורה אחת ותיבדק שבירת השורות מימין לשמאל.
\end{hebrew}
\end{document}
